\subsection{交换拓扑群的一致结构与完备化}

我们已经指出，在交换拓扑群 G 上左一致结构与右一致结构相同；每当我们谈到 G 的一致结构时，所指的就是这个唯一的一致结构。

定理 2 设 G 是一个交换拓扑群。则函数 \(x^{-1}\) 与 xy 分别在 G 与 G × G 上一致连续。此外 G 容许一个 Hausdorff 完备化 \(\hat{G}\)，且 \(\hat{G}\) 是交换的。

\(x^{-1}\) 的一致连续性由第 1 目命题 2 得出，而 xy 的一致连续性由第 1 目命题 3 得出，因为 \((x, y) \to xy\) 是 \(G \times G\) 到 G 中的一个连续同态。若 G 是 Hausdorff 的，它满足第 4 目定理 1 的条件（如同每个左一致结构与右一致结构相同的 Hausdorff 群一样）；此外由恒等延拓原理，函数 xy 与 yx 在 \(\hat{G} \times \hat{G}\) 上相等；由此得到定理的第二部分，办法是在一般情形考虑伴随于 G 的 Hausdorff 群。

特别地，由这个定理可得：若 \(f\) 与 \(g\) 是一致空间 \(\mathbf{X}\) 到一个交换群 \(G\)（以加法记的）中的两个一致连续映射，则函数 \(-f\) 与 \(f + g\) 是一致连续的。

命题 9 设 G 是一个交换群，\(\mathfrak{T}_1\)、\(\mathfrak{T}_2\) 是与 G 的群结构相容的两个 Hausdorff 拓扑。设 \(\mathfrak{T}_1\) 比 \(\mathfrak{T}_2\) 细，且存在关于 \(\mathfrak{T}_1\) 的 o 的一个基本邻域系，它的成员对 \(\mathfrak{T}_2\) 是闭的。设 \(\mathrm{G}_1, \mathrm{G}_2\) 分别是 G 关于拓扑 \(\mathfrak{T}_1, \mathfrak{T}_2\) 的完备化，\(f: \mathrm{G}_1 \to \mathrm{G}_2\) 是延拓 G 的恒等映射的连续同态（第 3 目命题 5）。则 \(f\) 是单射。

设 G 以加法记。设 \(\mathfrak{U}_{1}\) 是 G 上（依第 1 目）对应于拓扑 \(\mathfrak{T}_{1}\) 的一致结构：只需证明若 F 与 \(F'\) 是关于 \(\mathfrak{U}_{1}\) 的两个极小 Cauchy 滤子（第二章 § 3 第 2 目），它们在 \(G_{2}\) 中收敛于同一点 a，则 \(F = F'\)（第二章 § 3 第 7 目）。为此，只需证明 \(F \cap F'\) 是关于 \(\mathfrak{U}_{1}\) 的一个 Cauchy 滤子。设 V 是 G 中关于 \(\mathfrak{T}_{1}\) 的 o 的一个邻域，使得 V 在 \(\mathfrak{T}_{2}\) 中是闭的，又设 W 是 G 中关于 \(\mathfrak{T}_{1}\) 的 o 的一个对称邻域，使得 \(W + W \subset V\)。由假设，存在一个 \(W_{d}\)-小集 M（相应地 \(M'\)），它属于 F（相应地 \(F'\)）；若 \(x \in M\) 且 \(y \in M\)，我们有 \(y - x \in W\)，即 \(y \in x + W\)。若 \(\overline{W}\) 与 \(\overline{V}\) 是 W 与 V 在 \(G_{2}\) 中的闭包，则由此可得 \(y \in x + \overline{W}\)，因而由于 a 在 M 的闭包中，\(a \in x + \overline{W}\) 对每个 \(x \in M\) 成立。类似地，\(a \in x' + \overline{W}\) 对每个 \(x' \in M'\) 成立，于是 \(x - x' \in \overline{W} + \overline{W}\)；但由 \((x, y) \to x + y\) 是 \(G_{2} \times G_{2}\) 到 \(G_{2}\) 中的一个连续映射，我们有 \(\overline{W} + \overline{W} \subset \overline{W + W} \subset \overline{V}\)。由此可得，若 \(x \in M\) 且 \(x' \in M'\)，则 \(x - x' \in \overline{V} \cap G = V\)，因为 V 在 \(\mathfrak{T}_{2}\) 中是闭的；证毕。

推论 1 在命题 9 的假设下，若 A 是 G 的一个子集，它关于一致结构 \(\mathfrak{U}_{2}\)（对应于 \(\mathfrak{T}_{2}\) 的）是一个完备子空间，则 A 关于一致结构 \(\mathfrak{U}_{1}\)（对应于 \(\mathfrak{T}_{1}\) 的）也是一个完备子空间。

若 \(A_{1}\) 是 A 在 \(G_{1}\) 中的闭包，则 \(f(A_{1})\) 含于 A 在 \(G_{2}\) 中的闭包，而由假设这个闭包等于 A。由于依定义 \(f(A)=A\) 且 f 是单射，我们有 \(A_{1}=A\)。

推论 2 设 G 是一个交换群，\(\mathfrak{T}_1, \mathfrak{T}_2\) 是与 G 的群结构相容的两个拓扑。设 \(\mathfrak{T}_1\) 比 \(\mathfrak{T}_2\) 细，且存在 o 的一个基本邻域系 \(\mathfrak{B}\)（关于 \(\mathfrak{T}_1\) 的），它的成员关于一致结构 \(\mathfrak{U}_2\)（对应于 \(\mathfrak{T}_2\) 的）是完备的。则 G 关于一致结构 \(\mathfrak{U}_1\)（对应于 \(\mathfrak{T}_1\) 的）是完备的。

\(\mathfrak{B}\) 中的诸集在拓扑 \(\mathfrak{T}_{2}\) 中是闭的，因而由推论 1 关于一致结构 \(\mathfrak{U}_{1}\) 是完备的；于是结论由第 3 目命题 4 得出。

